\chapter{Memory Circuits: One-Loop, First-Order Systems (1-OS)}
% FIXME: \label{lect1}

\minitoc

In the second lesson we will introduce the memory circuits that are obtained through a first loop that brings to the input of a circuit the effect of applying a signal to another input. This reverse connection allows the circuit to maintain a state according to a command. Thus the memory function is obtained.

We will not go into details because we have a precise target: the internal structure of a computing system that includes a processor, with its internal storage resources, and the memories in which data and software are stored.

\section{Latch}

The simplest circuit will allow the latching of the simplest event: the temporary transition of a digital signal from 1 to 0, or the temporary transition of some circuit from 0 to 1.

\subsection{Closing the first loop}

Closing a feedback loop requires an output of a logic circuit to be connected to one of its inputs. In the Figure \ref{latches}a, the output of an AND gate is connected to one of the inputs, and the (Reset)' pulses are applied to the other input. If the transition to zero lasts long enough so that the signal propagates through the circuit to the second input, then the output is fixed at the zero value and the signal (Reset)' from the input can disappear. The transition event from the input is fixed to the output forever.

\placedrawing[H]{figures/03_01_latches.lp}{{\bf The elementary latches.}
{\small Using the  loop, closed from the output to one input, elementary storage elements are built. {\bf a}. {\em AND
loop} provides a {\em reset-only} latch. {\bf b}. {\em OR loop}
provides the {\em set-only} version of a storage element. {\bf c}.
The heterogeneous elementary {\em set-reset} latch results
combining the {\em reset-only} latch with the {\em set-only}
latch. {\bf d}. The wave forms describing the behavior of the previous three latch circuits.}}{latches}

Now let's take an OR gate and connect its output to one of the inputs, and on the other, apply a transition from 0 to 1. If the duration of the input signal allows the propagation of its effect on the reaction loop, the output of the circuit will be permanently fixed on the value 1.

The good news is that we can capture temporary events indefinitely. We say that the two circuits memorize, one the temporary transition to 0, the other the temporary transition to 1. The bad news is that we cannot make these circuits "forget" what they have memorized. But, by combining the two circuits, we manage to obtain a fundamental memory structure that can change its state according to the set command, S, to 1 or according to the reset command, R', to 0. The set command is active on 1, in while the reset command is active on 0.

\placedrawing[h]{figures/03_02_symmetric_latches.lp}{{\bf Symmetric elementary latches.}
{\small {\bf a.} Symmetric elementary {\em NAND latch} with
low-active commands S' and R'. {\bf b}. Symmetric elementary {\em
NOR latch} with high-active commands S and R.}}{symLatches}

We obtained a circuit with two states between which it can switch according to asymmetrical commands, one active on zero and the other active on one. For reasons of use, it is preferable that both commands are of the same type. We will use the two forms of Morgan's law to transform the circuit in Figure \ref{latches}c and we will obtain two symmetrical structures ordered with signals of the same type. Using the de Morgan rule  A + B = (A'B')', the latch from Figure \ref{latches}c becomes the latch from Figure \ref{symLatches}a, and using the de Morgan rule AB = (A'+B')', the latch from Figure \ref{latches}c becomes the latch from Figure \ref{symLatches}b.


\subsection{Clocked latch}

For the clocked latch two NAND gate are added (see Figure \ref{ck_latch}) to allow the R and R inputs conditioned by clock>

\placedrawing[H]{figures/03_03_clocked_latches.lp}{{\bf Elementary clocked latch.}
{\small The transparent RS clocked latch is sensitive
(transparent) to the input signals during the active level of the
clock (the high level in this example). {\bf a}. The internal
structure. {\bf b}. The logic symbol.}}{ck_latch}

{\bf The first latch problem}: the inputs for indicating
{\em how} the latch switches are the same as the inputs for
indicating {\em when} the latch switches; we must find a solution
for declutching the two actions   building a version  with
distinct inputs for specifying ``how" and ``when".

{\bf The second latch problem}: if we apply synchronously
S'=0 and R'=0 on the inputs of NAND latch  (or S=1 and R=1 on the
inputs of OR latch), i.e., the latch is commanded ``to switch in
both states simultaneously", then we can not predict what is the
state of the latch after the ending of these two active signals.

The first problem is solved in Figure \ref{ck_latch}a by conditioning the application of signals S' and R' to the latch in Figure \ref{symLatches}a by an additional signal that we will call clock, CK. Thus, we will use the S and R inputs to specify how the circuit switches, and the CK input to specify when the circuit switches. In this way, the decoupling of "{\em when}" from "{\em how}" is obtained.

For the second problem, we have a more "brutal" solution achieved through a restriction. We will connect between the two inputs a NOT that will not allow the simultaneous application of setting and resetting the circuit (see Figure \ref{d_latch}a). The problem is not solved. We avoid doing a wrong use.



\placedrawing[h]{figures/03_04_data_latch.lp}{{\bf The data latch.} {\small
Imposing the restriction $R=S'$ to an RS latch results the {\bf D
latch} without non-predictable transitions ($R=S=1$ is not anymore
possible). {\bf a.} The structure. {\bf b.} The logic symbol. {\bf
c.} An improved version for the data latch internal
structure.}}{d_latch}

For the time interval in which the clock is active, CK = 1, it says that the latch is transparent to the control signal S and R. The latch becomes non-transparent in the time interval in which CK = 0. Indeed, in the period in which the latch is transparent, the decoupling between "when" and "how" is not achieved, and the circuit switches conditioned by the evolution of the S and R inputs. A correct use of the latch assumes that during the transparency period the S and R inputs are stable.

We have to admit that we did not separate the "when" from the "how" rigorously enough. We will do it in the following.

\section{Serial extension: Master-Slave Principle}

The clocked latch, in the SR version as well as in the D version (data latch), allows switching at any time during the transparency. This fact limits the applications of this circuit. Many applications require switching precisely determined by the active, positive or negative transition of the clock signal. Applying the {\em master-slave principle} will allow this behavior.

Figure \ref{m_s}a shows a structure where transparency is blocked by the fact that the two RS latches have the clock applied in antiphase. On the active level of the clock, the positive one, the first latch, the master, is transparent, a fact that allows it to be switched according to the S and R inputs without affecting the output of the circuit because the second latch is not transparent. On the inactive level of the clock, the second latch, the slave, copies the state of the master which can no longer be changed. In this way, the entire circuit switches as a consequence of the negative transition of the clock signal. So, the active edge of the clock applied to the master-slave structure is the negative one. In Figure \ref{m_s}b, the logical symbol of the structure in Figure \ref{m_s}a is represented.

\placedrawing[H]{figures/03_05_master_slave.lp}{{\bf The master-slave principle.}
{\small Serially connecting two RS latches, activated with
different levels of the clock signal, results a non-transparent
storage element. {\bf a.} The structure of a RS master-slave
flip-flop, active on the falling edge of the clock signal. {\bf
b.} The logic symbol of the RS flip-flop triggered by the {\em negative edge} of clock. {\bf
c.} The logic symbol of the RS flip-flop triggered by the {\em positive edge} of clock.}}{m_s}

\placedrawing[H]{figures/03_06_d_flip_flop.lp}{{\bf The delay (D) flip-flop.}
{\small  Restricting the two inputs of an RS flip-flop to
$D=S=R'$, results an FF with predictable transitions. {\bf a.} The
structure. {\bf b.} The logic symbol. {\bf c.} The wave forms
proving the delay effect of the D flip-flop.}}{d_ff}


The second latch problem also propagates at the level of the master-slave structure. The solution we can come up with at this level is the one that introduces the limitation by connecting the inputs through an inverter circuit (see Figure \ref{d_ff}a). The result is the D-FF (delay flip-flop) circuit. The name is justified by the waveforms represented in Figure \ref{d_ff}c where the output of the circuit follows the evolution of the input with a delay equal to the period of the clock signal.

\section{Serial-parallel extension: Register}

\subsection{Structure}

The serial-parallel extension in the field of first-order systems (1-OS) is represented by the register. By connecting in parallel $n$ serially extended  structures (the delay flip-flop type master-slave circuit) is obtained a {\bf {\em register}} by applying the same clock signal on each D-FF (see Figure \ref{reg}a). A register stores an $n$-bit configuration synchronously with the active edge of the clock.

\placedrawing[h]{figures/03_07_register.lp}{{\bf The $n$-bit register.} {\small
{\bf a.} The structure: a bunch of DF-F connected in parallel.
{\bf b.} The logic symbol.} }{reg}

The Figure \ref{register} shows the effect of the transfer through a register.

\placedrawing[h]{figures/03_08_register_operation.lp}{{\bf Register operation.} {\small At each active
edge of clock (in this example it is the positive edge) the register's
output takes the value applied on its inputs if {\tt reset = 0} and
{\tt enable = 1}.}}{register}

The time behavior is mainly characterized by two times interval:
\begin{description}
    \item[minimal set-up time]: $t_{SU_{min}}$ is the minimum time interval that the input must be kept stable before the active edge of the clock so that the transition of the output corresponds to the input
    \item[maximal propagation time]: $t_p$ is the propagation time to the output related to the active edge of clock.
\end{description}

\subsection{Applications}

Follows the list of the main applications of the register in digital systems.

\subsubsection{Storing}
The {\tt enable} input allows us to determine when (i.e., in what
clock cycle) the input is loaded into a register. If {\tt enable =
0}, the registers {\em stores} the data loaded in the last clock
cycle when the condition {\tt enable = 1} was fulfilled. This
means we can keep the content once stored into the register as
much time as it is needed.

\subsubsection{Buffering}
The registers can be used to {\em buffer} (to isolate, to
separate) two distinct blocks so as some behaviors are not
transmitted through the register. For example, in Figure
\ref{register} the transitions from {\tt c} to {\tt d} and from
{\tt d} to {\tt e} at the input of the register are not
transmitted to the output.

\subsubsection{Synchronizing}
For various reasons the digital signals are generated ``unaligned
in time" to the inputs of a system, but they are needed to be
received very well controlled in time. We say usually, the signals
are applied {\em asynchronously} but they must be received {\em
synchronously}.  For example, in Figure \ref{register} the input
of the register changes somehow chaotically related to the active
edge of the clock, but the output of the register switches with a
constant delay after the positive edge of clock. We say the inputs
are synchronized to the output of the register. Their behavior is
``time tempered".

\subsubsection{Delaying}
The input value applied in the clock cycle {\tt n} to the input of
a register is generated to the output of the register in the clock
cycle {\tt n+1}. In other words, the input of a register is
delayed one clock cycle to its output. See in Figure
\ref{register} how the occurrence of a value in one clock cycle to
the register's input is followed in the next clock cycle by the
occurrence of the same value to the register's output.

\subsubsection{Looping}
Structuring a digital system means to make different kind of
connections. One of the most special, {\em as we see in what
follows}, is a connection from some outputs to certain inputs in a
digital subsystem. This kind of connections are called {\bf
loops}. The register is an important structural element in closing
controllable loops inside a complex system.

\subsubsection{Pipelining}

\begin{exemplu}
The previously exemplified {\tt threeAdder.sv} module performs the addition of three numbers in twice the time of the sum of two numbers. We can reduce the execution time through a pipeline solution that uses two registers. The solution is presented in the following module:
{\small
{\em
\begin{shaded*}
\begin{lstlisting}[language=Verilog, caption= , morekeywords={always_ff, always_comb, logic, loop, repeat, doInParallel, for}]
/*************************************************************************
File name:      pipelinedThreeAdder.sv
Circuit name:   pipelined Three Input Adder
Description:    The module 'threeAdder' has 3 4-bit inputs and one 4-bit
                output. The circuit adds modulo 16 three numbers;
                do not provide carry output
*************************************************************************/
module pipelinedThreeAdder( output  logic [3:0] out     ,
                            input   logic [3:0] in0     ,
                            input   logic [3:0] in1     ,
                            input   logic [3:0] in2     ,
                            input   logic       clock   );
    logic [3:0] syncReg ;
    logic [3:0] pipeReg ;
    logic [3:0] sum     ;

    adder   inAdder(    .out(sum    ),
                        .in0(in0    ),
                        .in1(in1    )),
            outAdder(   .out(out    ),
                        .in0(syncReg),
                        .in1(pipeReg));
    always_ff @(posedge clock)  begin   syncReg <= in2  ;
                                        pipeReg <= sum  ;
                                end
endmodule
\end{lstlisting}
\end{shaded*}
}
}

Two modules of {\tt adder} defined Example \ref{adder}
are instantiated as {\tt inAdder} and {\tt outAdder}, they are
interconnected using the pipeline  register {\tt pipeReg}. The register {\tt subcReg} is used to synchronize the {\tt in2} input with the pipelined result provided by {\tt inAdder}.

$\diamond$
\end{exemplu}


\section{Parallel extension: Random-Access Memory}

The parallel extension in first-order systems (1-OS) is represented by random access memory (RAM).

\subsection{Generic structure}

The generic structure of a memory with $2^p$ one-bit locations is represented in the Figure \ref{ram}.

\placedrawing[h]{figures/03_09_ram.lp}{{\bf The principle of the random
access memory (RAM).} {\small The clock is distributed by a DMUX
to one of $m=2^{p}$ DLs, and the data is selected by a MUX from
one of the $m$ DLs. Both, DMUX and MUX use as selection code a
$p$-bit address. The one-bit data DIN can be stored in the clocked
DL.}}{ram}

\placedrawing[H]{figures/03_10_async_ram.lp}{{\bf Read and write cycles for an
asynchronous RAM.} {\small Reading is a combinational process of
selecting. The access time, $t_{ACC}$, is given by the propagation
through a big MUX. The write enable signal must be strictly
included in the time interval when the address is stable (see
$t_{ASU}$ and $t_{AH}$). Data must be stable related to the
positive transition of $WE'$ (see $t_{DSU}$ and $t_{DH}$).}}{wave}


The waveforms that define the operation of the structure in Figure \ref{ram} are represented in Figure \ref{wave} where the main restrictions are represented by:
\begin{description}
    \item[$t_{ASU}$]: address set-up time represents the time interval in which the address must be stable before the activation of the WE signal to allow the decoder in the DMUX to do its work.
    \item[$t_{DSU}$]: data set-up time represents the time interval in which DIN must be stable before de-activating the WE signal to allow the correct closing of the reaction loop in the selected latch.
    \item[$t_{DH}$]: data hold represents the time interval in which DIN must be stable after deactivating the WE signal to allow the correct closing of the reaction loop in the selected latch.
    \item[$t_{W}$]: width of the write enable signal which ensures correct writing in the selected latch
    \item[$t_{AH}$]: address hold time is the time interval in which the address must be maintained after the WE signal is provided
    \item[$t_{ACC}$]: access time is the time interval after which the content of the selected cell is accessible at the output
\end{description}
All the previously listed times are minimum values.

\placedrawing[H]{figures/03_11_async_mbit_ram.lp}{{\bf The asynchronous $m$-bit word
RAM.} {\small Expanding the number of bits per word means to
connect in parallel one-bit word memories which share the same
decoder. Each COLUMN contains storing latches and  AND-OR
circuits for one bit.}}{bigword}\textcolor[rgb]{0.00,0.00,0.00}{}


How can you build a RAM memory that stores words of $m$ bits? The generic solution is presented in Figure \ref{bigword} where for each bit a column containing $2^p$ latches is defined.

In Figure \ref{bigword}, $DCD_p$ is shared between the DMUX and the MUX in Figure \ref{ram}. The ANDs associated with the demultiplexing function are selected with the WE signal. The AND-OR structure is repeated $m$ times, once in each $COLUMN_i$ column.

\subsection{Synchronous RAM}\label{syncMem}

At very high speeds, on the order of GHz, the time restrictions illustrated in Figure \ref{wave} are very difficult to fulfill in the already described asynchronous version. To increase the degree of controllability of our projects, synchronous RAM memories are used. In Figure \ref{swave}, the behavior of a synchronous memory (SRAM) is illustrated in the sense that all the time intervals that characterize the behavior of the memory are related to the active edge of a clock signal. In this case, the designer of a system that includes a synchronous memory can more rigorously control the time behavior of the system.

For the memories used by the system designer, {\em memory libraries} are provided by specialized companies in the form of IPs (intellectual property), guaranteeing the correctness of the code used in the description.

\placedrawing[h]{figures/03_12_sync_ram.lp}{{\bf Read and write cycles for  Synchronous RAM (SRAM).}
{\small For the {\em flow-through} version of a $SRAM$ the time
behavior is similar to a register. The set-up and hold time are
defined related to the active edge of clock for all the input
connections: {\em data}, {\em write-enable}, and {\em address}.
The data output is also related to the same edge.}}{swave}

{\small
\begin{shaded*}
\begin{lstlisting}[language=Verilog, caption= , morekeywords={always_ff, always_comb, logic, loop, repeat, doInParallel, for}]
/*************************************************************************
File name:      syncRAM.sv
Circuit name:   Synchronous RAM
Description:    4096 32-bit words synchronous RAM; asynchronous read
*************************************************************************/
module syncRAM( output  logic [31:0]    out     ,
                input   logic [31:0]    in      ,
                input   logic [12:0]    addr    ,
                input   logic           we      ,
                input   logic           clock   );
    logic [31:0]    mem[0:8191]   ;

    always_ff @(posedge clock) if (we) mem[addr] <= in  ;
    assign out  = mem[addr] ;
endmodule
\end{lstlisting}
\end{shaded*}
}

\subsection{Synchronous pipelined RAM}\label{sybcPipeMem}

The time $t_{ACC}$ offered by a RAM can sometimes be too high for the speed requirements of the system. To increase the clock frequency, a register is used at the memory output. Thus, the data is obtained at the new output very quickly: $t_{ACC}$ is the propagation time through the register. But there is a price: the {\bf {\em latency}} of one clock cycle (we remember that the register consists of {\em delay} flip-flops).

A skilled designer will almost always be able to hide the effect of latency and take advantage of the increase in system frequency obtained through the pipeline register from the RAM output.

{\small
\begin{shaded*}
\begin{lstlisting}[language=Verilog, caption= , morekeywords={always_ff, always_comb, logic, loop, repeat, doInParallel, for}]
/*************************************************************************
File name:      syncPipeRAM.sv
Circuit name:   Synchronous pipelined RAM
Description:    4096 32-bit words synchronous pipelined RAM
*************************************************************************/
module syncPipeRAM( output  logic [31:0]    out     ,
                    input   logic [31:0]    in      ,
                    input   logic [12:0]    addr    ,
                    input   logic           we      ,
                    input   logic           clock   );
    logic [31:0]    mem[0:8191]   ;

    always_ff @(posedge clock)  begin if (we)   mem[addr] <= in     ;
                                                out  <= mem[addr]   ;
                                end
endmodule
\end{lstlisting}
\end{shaded*}
}

\subsection{{\bf {\em Register file}}}

A register file is a battery of registers from which, in each clock cycle, any two registers can be accessed for reading and any register for writing. The implementation of such a circuit is done with a synchronous memory with three ports: two for reading and one for writing (see Figure \ref{file_reg}), where:

\placedrawing[h]{figures/03_13_reg_file.lp}{{\bf Register file.} {\small In this
example it contains $2^{n}$ $m$-bit registers. In each clock cycle
any two registers can be read and writing can be performed in
anyone. }}{file_reg}

\begin{description}
    \item{{\tt leftAddr[n-1:0]}}: is the address used to select data to the {\tt leftOut[m-1:0]} output
    \item{{\tt rightAddr[n-1:0]}}: is the address used to select data to the {\tt rightOut[m-1:0]} output
    \item{{\tt destAddr[n-1:0]}}: is the address used to select the location where the value applied on input {\tt in[m-1:0]} is stored at the positive edge of the {\tt clock} signal
    \item{{\tt writeEnable}}: is the write enable command.
\end{description}
The internal structure of a file register is characterized by the two multiplexers that ensure access to any pair of stored values.

{\small
\begin{shaded*}
\begin{lstlisting}[language=Verilog, caption= , morekeywords={always_ff, always_comb, logic, loop, repeat, doInParallel, for}]
/*************************************************************************
File name:      regFile.sv
Circuit name:   Register File
Description:    Three-port, 32 32-bit words implemented with a
                synchronous RAM
*************************************************************************/
module regFile( output  logic [31:0]    leftOut     ,
                output  logic [31:0]    rightOut    ,
                input   logic [31:0]    in          ,
                input   logic [4:0]     leftAddr    ,
                input   logic [4:0]     rightAddr   ,
                input   logic [4:0]     destAddr    ,
                input   logic           we          ,
                input   logic           clock       );
    logic [31:0]    mem[0:31]   ;

    always_ff @(posedge clock) if (we) mem[destAddr] <= in  ;

    assign leftOut  = mem[leftAddr] ;
    assign rightOut = mem[rightAddr];
endmodule
\end{lstlisting}
\end{shaded*}
}

\begin{verisum}{\em :

\hspace{-0.8cm} \rule[5mm]{16cm}{0.25mm} \vspace{-0.7cm}

\begin{description}
    \item[logic] [m-1:0] mem[0:n-1]: defines a block of memory of $n$ $m$-bit words
    \item[always\_ff]: defines a block which describes the behavior of register-like sequential circuits; it is followed always at least by {\tt {\bf @(posedge} clock)} or {\tt {\bf @(negedge} clock)} specifying the clock signal and its active edge. The non-blocking assignment, {\tt <=}, is used to assure the sequential behavior. While {\tt always\_comb} describes the behavior of a combinational circuit, {\tt always\_ff} ({\tt ff} comes from flip-flop) describes the behavior of the sequential, clock triggered, circuit called {\tt register}. Now {\tt logic} represents a register instead of an output of a combinational circuit like in the case of {\tt always\_comb}.
\end{description}
\rule[5mm]{16cm}{0.25mm} }
\end{verisum}


%****************************************************************************
\subsection{Representation of \( B_n \) in MSB and LSB Forms with Endianness}
%****************************************************************************
The set \( B_n \), representing \( n \)-bit unsigned binary numbers, can be expressed using two common conventions for bit ordering: Most Significant Bit (MSB) first and Least Significant Bit (LSB) first. These conventions are closely tied to the concept of endianness in computer architecture, which determines how binary numbers are stored or transmitted in memory.

%****************************************************
\subsubsection{MSB Representation (Big-Endian Order)}
%****************************************************
In the Most Significant Bit (MSB) representation, the leftmost bit (\( b_{n-1} \)) corresponds to the highest positional value (\( 2^{n-1} \)), and the rightmost bit (\( b_0 \)) corresponds to the lowest positional value (\( 2^0 \)). This is the standard positional notation for binary numbers. 

For \( B_n \), the value of \( x \) is calculated as:
\[
x = b_{n-1} \cdot 2^{n-1} + b_{n-2} \cdot 2^{n-2} + \dots + b_1 \cdot 2^1 + b_0 \cdot 2^0, \quad b_i \in \{0, 1\}.
\]

%*******************************************************
\subsubsection{LSB Representation (Little-Endian Order)}
%*******************************************************
In the Least Significant Bit (LSB) representation, the least significant bit (\( b_0 \)) is stored or transmitted first. However, the numerical value of the number remains the same, as the binary digits are still interpreted using the same positional notation. The formula remains:
\[
x = b_{n-1} \cdot 2^{n-1} + b_{n-2} \cdot 2^{n-2} + \dots + b_1 \cdot 2^1 + b_0 \cdot 2^0, \quad b_i \in \{0, 1\}.
\]

The difference lies in the physical storage or transmission order:
- In big-endian systems, bits are stored or transmitted starting with the most significant bit (\( b_{n-1} \)).
- In little-endian systems, bits are stored or transmitted starting with the least significant bit (\( b_0 \)).


Example for \( n = 3 \)

For \( B_3 = \{0, 1, 2, 3, 4, 5, 6, 7\} \), the logical representation (big-endian) is:

\textbf{Logical Representation (Big-Endian)}:
\[
\begin{aligned}
x = 0 & \Rightarrow 000, \\
x = 1 & \Rightarrow 001, \\
x = 2 & \Rightarrow 010, \\
x = 3 & \Rightarrow 011, \\
x = 4 & \Rightarrow 100, \\
x = 5 & \Rightarrow 101, \\
x = 6 & \Rightarrow 110, \\
x = 7 & \Rightarrow 111.
\end{aligned}
\]

In a little-endian system, the storage or transmission order of bits is reversed. For instance:
- \( x = 1 \) (big-endian: \( 001 \)) is stored as \( 100 \) (little-endian).
- \( x = 6 \) (big-endian: \( 110 \)) is stored as \( 011 \) (little-endian).

\textbf{Physical Representation in Little-Endian (Stored Order)}:
\[
\begin{aligned}
x = 0 & \Rightarrow 000, \\
x = 1 & \Rightarrow 100, \\
x = 2 & \Rightarrow 010, \\
x = 3 & \Rightarrow 110, \\
x = 4 & \Rightarrow 001, \\
x = 5 & \Rightarrow 101, \\
x = 6 & \Rightarrow 011, \\
x = 7 & \Rightarrow 111.
\end{aligned}
\]

%***************************************
\subsubsection{Commentary on Endianness}
%***************************************
- **Big-Endian Systems**:
  - Big-endian systems store numbers in memory starting with the most significant bit. This is often the default in higher-level representations of binary numbers and is consistent with standard positional notation.

- **Little-Endian Systems**:
  - Little-endian systems store numbers in memory starting with the least significant bit. This convention simplifies certain hardware operations, such as incrementing binary counters, and is widely used in modern processors like x86 architectures.

- **Practical Implications**:
  - While the logical representation of binary numbers is the same regardless of endianness, understanding the storage or transmission order is crucial for tasks like data serialization, memory addressing, and communication protocols.

%%%%%%%%%%%%%%%%%%%%%%%%%%%%%%%%%%
%\subsection{Mathematical Preference and Use of MSB and LSB Expansions}
%%%%%%%%%%%%%%%%%%%%%%%%%%%%%%%%%%

The binary expansions of \( x \) differ in the ordering of bits depending on whether the **Most Significant Bit (MSB)** or **Least Significant Bit (LSB)** form is used. These differences arise from practical applications in engineering versus conventions in mathematics.

Mathematical Preference: MSB Form

Mathematicians universally use the **MSB expansion** to represent binary numbers. The MSB form is consistent with the positional arithmetic taught in grade school, where each digit corresponds to a specific positional value. For \( n \)-bit unsigned binary numbers, the MSB expansion is:

\[
x = b_{n-1} \cdot 2^{n-1} + b_{n-2} \cdot 2^{n-2} + \dots + b_1 \cdot 2^1 + b_0 \cdot 2^0, \quad b_i \in \{0, 1\}.
\]

This form aligns naturally with:
\begin{itemize}[noitemsep, left=1cm, topsep=0pt]
    \item **Grade school arithmetic**: Positional notation, where the leftmost digit has the highest weight, and each subsequent digit corresponds to decreasing powers of the base.
    \item **Decimal analogy**: For example, the decimal number \( 345 \) is interpreted as:
    \[
    345 = 3 \cdot 10^2 + 4 \cdot 10^1 + 5 \cdot 10^0.
    \]
    Similarly, in binary, \( 110 \) is interpreted as:
    \[
    110 = 1 \cdot 2^2 + 1 \cdot 2^1 + 0 \cdot 2^0.
    \]
\end{itemize}

Thus, the MSB form is preferred in mathematics for its clarity and compatibility with standard positional arithmetic.

---

Use of LSB Form in Engineering and Computing

In contrast, the **LSB expansion** reverses the storage or transmission order of bits, prioritizing the least significant bit (\( b_0 \)). The formula remains the same, but the physical ordering differs, reflecting how the bits are handled in memory or transmitted serially. For example:
\[
x = b_0 \cdot 2^0 + b_1 \cdot 2^1 + \dots + b_{n-2} \cdot 2^{n-2} + b_{n-1} \cdot 2^{n-1}.
\]

The LSB form is not used in pure mathematics but is important in:
\begin{itemize}[noitemsep, left=1cm, topsep=0pt]
    \item **Computer systems**:
        - Little-endian architectures (e.g., x86) use LSB-first order for storing and transmitting binary numbers.
        - This simplifies certain hardware operations, such as incrementing binary counters or performing addition bit-by-bit starting from the least significant bit.
    \item **Data communication**:
        - Serial communication protocols often transmit the least significant bit first due to hardware design constraints.
\end{itemize}

---

Commentary on MSB as Positional Arithmetic

The MSB expansion closely resembles the positional arithmetic taught in elementary education. Each digit (or bit) contributes to the total value based on its position, with weights corresponding to powers of the base (in this case, \( 2 \)). This positional system is universally understood and forms the foundation for number systems across different bases (e.g., decimal, binary, hexadecimal).

In contrast, the LSB form prioritizes computational efficiency in specific applications rather than logical clarity. While it has practical importance in engineering and computing, it is rarely used for manual or theoretical calculations.

---

Conclusion

Mathematicians exclusively use the MSB form due to its consistency with positional arithmetic and its alignment with grade school notation. The LSB form is primarily a practical convention in computer architecture and data communication, where hardware efficiency and specific operational needs dictate its use.



%%%%%%%%%%%%%%%%%%%%%%%%
\section{Problems}
%%%%%%%%%%%%%%%%%%%%%%%%
\subsection{Registers}

\begin{problem}\label{mean}
Design a system that receives a signed integer in each clock cycle and outputs the sum of the last four received numbers. When the system is reset, it is considered that the last four received numbers had the value 0.

If $t_{SU_{min}} = \#1$, $t_p = \#5$, $t_{sum} = \#50$, then what is the minimum period of the clock at which the system can present a correct result when entering a register.

$\diamond$
\end{problem}

{\bf Solution}

The system consists in 4 registers, {\tt reg0, ..., reg3}, serially connected and a tree of three adders used to sum the content of the registers. The first register, {\tt reg0}, receives in each clock cycle the input value, {\tt in}. Register {\tt reg1} receives the content of {\tt reg1}, {\tt reg2} receives the content of {\tt reg1}, and {\tt reg3} the content of {\tt reg2}. The first two and the last two registers are added using two adders whose outputs are added using the output adder. The division by 4 is simply performed selecting the 32 most significant bits from the output of the last adder.
{\small
\begin{shaded*}
\begin{lstlisting}[language=Verilog, caption= , morekeywords={always_ff, always_comb, logic, loop, repeat, doInParallel, for}]
/*************************************************************************
File name:      mean.sv
Circuit name:
Description:
*************************************************************************/
module mean(output  logic [31:0]    out     ,
            input   logic [31:0]    in     	,
            input   logic           reset  	,
            input   logic           clock   );
    logic [31:0]    reg0;
    logic [31:0]    reg1;
    logic [31:0]    reg2;
    logic [31:0]    reg3;
    logic [32:0]    sum0;
    logic [32:0]    sum1;
    logic [33:0]    sum;

    always_ff @(posedge clock)  if( reset)  begin   reg0 <= 0   ;
                                                    reg1 <= 0   ;
                                                    reg2 <= 0   ;
                                                    reg3 <= 0   ;
                                            end
                                    else    begin   reg0 <= in  ;
                                                    reg1 <= reg0;
                                                    reg2 <= reg1;
                                                    reg3 <= reg2;
                                            end
    assign #1 sum0 = reg0 + reg1    ;
    assign #1 sum1 = reg2 + reg3    ;
    assign #1 sum  = sum0 + sum1    ;
    assign    out  = sum[33:2]      ;
endmodule

module testMean;
    logic [31:0]    out     ;
    logic [31:0]    in     	;
    logic           reset  	;
    logic           clock   ;

    initial begin               clock = 0       ;
                    forever #3  clock = ~clock  ;
            end

    initial begin       reset = 1       ;
                        in    = 32'b1000;
                    #6  reset = 0       ;
                    #60 $finish         ;
            end

    mean dut(   out     ,
                in      ,
                reset   ,
                clock   );
endmodule
\end{lstlisting}
\end{shaded*}
}
The minimal period of clock is: $T_{clock} = t_p + 2 \times t_{sum} + t_{SU_{min}} = \#106$

\begin{problem}
Revisit the Problem \ref{mean} and insert the first two adders and the output adder pipe registers.
\begin{enumerate}
    \item What is the resulting minimal period of clock
    \item Design this pipelined version
    \item Simulate the resulting design.
\end{enumerate}

$\diamond$
\end{problem}

\begin{problem}
Design a fully buffered ALU and evaluate the minimal period of the clock signal if $t_{SU_{min}} = \#1$, $t_p = \#5$, $t_{ALU} = \#60$. Fully buffered means inputs on ALU are provided from registers and the output is send through a register.

$\diamond$
\end{problem}

\begin{problem}
Design a system that receives a string of bytes synchronized with the clock and transfers it synchronously with a selectable delay. The selection is made with {\tt sel[1:0]} indicating a delay of 1, 2, 3, or 4 clock cycles.

$\diamond$
\end{problem}

\subsection{Memories}

\begin{problem}
Design a synchronous pipelined RAM for $2^{addressSize}$ $(wordSize)$-bit words.

$\diamond$
\end{problem}

{\bf Solution}

{\small
\begin{shaded*}
\begin{lstlisting}[language=Verilog, caption= , morekeywords={always_ff, always_comb, logic, loop, repeat, doInParallel, for}]
/*************************************************************************
File name:      defines.vh
*************************************************************************/
`define wordSize     32
`define addressSize  16
\end{lstlisting}
\end{shaded*}
}

{\small
\begin{shaded*}
\begin{lstlisting}[language=Verilog, caption= , morekeywords={always_ff, always_comb, logic, loop, repeat, doInParallel, for}]
/*************************************************************************
File name:      syncPipeRAM.sv
Circuit name:   Synchronous pipelined RAM
Description:    2^(`addressSize) locations of  (`wordSize)-bit words
                synchronous pipelined RAM
*************************************************************************/
`include "DEFINES.vh"
module syncPipeRAM( output  logic [`wordSize-1:0]       out     ,
                    input   logic [`wordSize-1:0]       in      ,
                    input   logic [`addressSize-1:0]    addr    ,
                    input   logic                       we      ,
                    input   logic                       clock   );
    logic [`wordSize-1:0]    mem[0:2**{`addressSize}-1]   ;

    always_ff @(posedge clock)  begin if (we)   mem[addr] <= in     ;
                                                out  <= mem[addr]   ;
                                end
endmodule
\end{lstlisting}
\end{shaded*}
}

\begin{problem}
Consider memory modules of 1KB and design with them a memory of 4K 32-bit word.

$\diamond$
\end{problem}

\begin{problem}
Simulate the behavior of the register file described in Section \ref{fastregFile} to prove the read-during-write operation.

$\diamond$
\end{problem} 